\begin{abstract}
En la operación de comercio electrónico de un supermercado, la demanda de pedidos online presenta marcadas fluctuaciones horarias y semanales que dificultan dimensionar la dotación de personal en el sector de preparación.
Una capacidad insuficiente frente a los picos de arribos genera acumulación de órdenes en fila y demoras en la entrega, mientras que el sobredimensionamiento incrementa la capacidad ociosa.
Este trabajo desarrolla y evalúa un modelo cuantitativo de apoyo a la decisión para determinar cuántos preparadores asignar en escenarios de demanda baja, media y alta con el objetivo de que al menos el 90~\% de los pedidos queden listos en un plazo máximo de 30~minutos desde su ingreso.
La metodología encadena tres técnicas de Investigación Operativa sobre un histórico sintético reproducible de 26~semanas (35\,402~pedidos) calibrado a partir del sistema ERP desarrollado por el equipo: primero, se pronostica el perfil horario de llegadas evaluando métodos estacionales sobre un período de prueba reservado, donde la media estacional de ocho semanas reduce el error absoluto medio un 24~\% respecto del método ingenuo estacional (de 4,51 a 3,42~pedidos por hora); segundo, se formula el sistema de fila común y se verifica el motor computacional contra las expresiones analíticas de régimen estacionario del modelo $M/M/c$ y la ley de Little; tercero, se ejecuta una simulación de eventos discretos en SimPy con números aleatorios comunes (100~réplicas por escenario y dotaciones de 2 a 8~preparadores) que incorpora tasas horarias variables y arrastre de pendientes entre horas.
Los resultados determinan una dotación mínima recomendada de $c^* = 3$~preparadores para jornadas de demanda baja (domingo, 92,0~\% de cumplimiento), $c^* = 4$ para demanda media (jueves, 97,4~\%) y $c^* = 6$ para demanda alta (sábado, 97,5~\%).
El contraste sobre los pedidos observados en las cuatro semanas de prueba confirma que dichas dotaciones sostienen el objetivo operativo (90,6~\%, 100,0~\% y 98,1~\% de pedidos a tiempo), y el análisis de sensibilidad cuantifica los márgenes de absorción ante subestimaciones del pronóstico y variaciones en los tiempos unitarios de armado.
\end{abstract}
